\documentclass[../../main.tex]{subfiles}
\begin{document}

\begin{flushright}
\footnotesize\textit{【自動文字起こし・要確認】}
\end{flushright}

{\bf [解]}
\paragraph{(1)}

時刻 $t=0$ に穴をとび出した水の，時刻 $t$ でのイチを $(X,Y)$ とする．\\
ただし，右のように $XY$ 平面をとる．

\begin{figure}[htb]
\centering
\begin{tikzpicture}[scale=1.0, >=stealth]
  \draw[->] (-0.5,0) -- (4.0,0) node[right] {$x$};
  \draw[->] (0,-0.5) -- (0,3.0) node[above] {$y$};
  \node[below left] at (0,0) {$O$};

  \node[left] at (0,1.8) {$a-h$};
  \fill (0,2.3) circle (1.5pt) node[left] {噴出口};
  \draw[->,thick] (0,1.8) -- (0.8,1.8);
  \draw[domain=0:3.2,samples=100,smooth,thick] plot (\x, {1.8 - 0.15*(\x)^2});
\end{tikzpicture}
\caption{噴出口から放物線状に落下する水の軌道}
\end{figure}

\begin{align}
\begin{cases}
X = \sqrt{2gh} \, t \\
Y = a - h - \dfrac{1}{2}gt^2
\end{cases}
\end{align}
$\sqrt{2gh} \neq 0$ から，$t$ をけして
\begin{align}
Y &= -\dfrac{1}{2}g \cdot \dfrac{X^2}{2gh} + a - h \\
&= -\dfrac{1}{4h} X^2 - h + a \quad \cdots \text{①}
\end{align}
したがって，水平到達キョリは $X = 2\sqrt{h(a-h)}$

\paragraph{(2)}

(1)の値を最大にするのは $h = \dfrac{1}{2}a$，つまり水槽の中点．

\paragraph{(3)}

①で $0 < h < a$ とするときの $(X,Y)$ のとる値を考える．
\begin{align}
Y' = -1 + \dfrac{X^2}{4h^2} =
\end{align}
だから，下表をとる． ($\because$ (2)から $0 \leqq X < a$)

\begin{table}[htb]
\centering
\caption{$Y$ の増減表}
\begin{tabular}{c|ccccc}
$h$ & $0$ & $\dots$ & $\dfrac{X}{2}$ & $\dots$ & $a$ \\
\hline
$Y'$ & & $+$ & $0$ & $-$ & \\
\hline
$Y$ & $(-\infty)$ & $\nearrow$ & $a-X$ & $\searrow$ & $-\dfrac{X^2}{4a}$
\end{tabular}
\end{table}

これと $0 < Y$ から，
\begin{align}
0 < Y \leqq a - X
\end{align}
が求める領域である．

\end{document}
